\section{Registro delle decisioni supportate da AI}

Sono riportate quattro interazioni che hanno influito sul sistema.
Le richieste sono sintetizzate; le motivazioni esplicitano il criterio
tecnico delle scelte adottate.

\begin{enumerate}
    \item \textbf{Modello del costo rimanente.}
    Richiesta: come addestrare un predittore a partire dalle feature e
    dai costi corretti. Proposta: usare una regressione lineare con
    \texttt{scikit-learn}, addestrata con \texttt{fit}.
    Esito: accettata e implementata.
    Motivazione: il modello consente di esaminare esplicitamente i quattro
    pesi e l'intercetta e di salvare una formula di previsione semplice.

    \item \textbf{Separazione dei dati per istanza.}
    Richiesta: come preparare i dati per l'addestramento e la verifica.
    Proposta: conservare l'identificativo di ogni istanza e separare gli
    otto seed di addestramento dai due seed di test.
    Esito: accettata e implementata.
    Motivazione: stati dello stesso problema condividono le coordinate;
    distribuirli tra addestramento e test renderebbe meno significativo
    il controllo su istanze non usate per apprendere i pesi.

    \item \textbf{Uso della previsione nella ricerca.}
    Richiesta: come passare dalla previsione di un singolo stato a una
    ricerca predittiva. Proposta: ordinare i successori tramite
    $g+\widehat h$, inserendoli in ordine decrescente nella pila LIFO,
    e mantenere la potatura sul solo costo accumulato.
    Esito: accettata e implementata in \texttt{predictiveSearch}.
    Motivazione: esplorare prima le stime più basse senza trattare
    una previsione appresa come limite inferiore certificato.

    \item \textbf{Valutazione della qualità a budget limitato.}
    Richiesta: come interpretare il costo trovato quando la ricerca
    non è completa. Proposta: calcolare un ottimo di riferimento con
    una ricerca separata e riportare l'\emph{optimality gap}.
    Esito: accettata e implementata; il punto di ingresso è stato
    riorganizzato in generazione, addestramento e confronto.
    Motivazione: distinguere la qualità del giro trovato dalla
    certificazione di ottimalità e confrontare i metodi sullo stesso
    budget online, senza usare l'ottimo per guidarli.
\end{enumerate}
